\subsection{A prescribed bar boundary can fail in an admissible theory}

\begin{proposition}[C31: \status{Proved}; failure of automatic bar extension]
\label{sa08:C31}
There is an admissible generalized theory $A$ over $k=\Z_{(2)}$,
with flat terms and $A^{w,q}=0$ for $q>w$, and a vertically closed
simplicial map
\[
 P=K(k,2)\longrightarrow A^{1,1},
\]
such that its map on the original odd disk of $U_1(P)$ does not
extend to a strict simplicial PD-DGA map $U_1(P)\to A$.
The failure occurs already in arity two. Equivalently, after adjoining
the coefficient algebra, the prescribed map does not extend from the
original disk to $S\otimes_k U_1(P)$ with the coefficients fixed.
\end{proposition}
\begin{proof}
Form the even-to-odd PD Poincar\'e algebra
\[
 D(P)=\Gamma_k(Pa)\otimes_k\Lambda_k(Px),\qquad
 da=x,\quad dx=0,
\]
where $a$ has form degree zero, $x$ has form degree one,
and both have weight one. Give both parameter arity one and let
$V$ be the quotient of $D(P)$ by all parameter arities at least three.
This is a differential PD ideal: the differential preserves arity,
products add it, and higher divided powers multiply it. Thus $V$
has only parameter weights zero, one and two.

At every simplex the positive fixed-arity Poincar\'e complexes split
universally. For one basis variable, the positive-arity complex is
\[
 k\gamma_m(a)\xrightarrow{\ 1\ }
 k\gamma_{m-1}(a)x.
\]
For any basis multidegree, contract its first nonzero tensor factor.
These are integral splittings and remain splittings after arbitrary
base change. Hence $V\to k$ is a universal vertical quasi-isomorphism,
and its terms are free. No simplicially natural contraction is needed.

Let $G_4$ be the elementary source of Statement~\ref{sa08:C09}
with marked weight four, which is defined because four is a power
of two. Put
\[
 A=G_4\otimes_k V,
\]
with coefficients $S=\Gamma_k(T)$ supplied by $G_4$ and the tensor
augmentation. The Poincar\'e splittings of the two factors show that
$A$ resolves $S$ vertically in every weight and every simplex,
universally under coefficient base change. Its terms are flat, and
the support condition $q\leq w$ holds.

For completeness, a concrete horizontal threshold through form degree
$d$ is
\[
 W\geq4(d+1)+2.
\]
Indeed every summand of $A^{W,q}$ is a tensor product of a $V$ row
of weight $t\in\{0,1,2\}$ and form degree $j$ with
$G_4^{W-t,q-j}$. The displayed inequality places this latter row
in its acyclic range $W-t\geq4(q-j+1)$. Its flatness and
Eilenberg--Zilber imply acyclicity after tensoring with the $V$ row.
There are only finitely many summands. This verifies uniform eventual
horizontal acyclicity, so $A$ is admissible.

In weights less than four, $G_4$ has only constant coefficient
terms in form degree zero. Therefore
\[
 A^{2,2}=\Lambda^2P.
\]
The map $P\to A^{1,1}$ sending the parameter to $x$ is simplicial
and vertically closed. It defines the prescribed map on the original
odd disk of $U_1(P)$ by sending its even differential $y$ to zero.

The arity-two, form-degree-two term of $U_1(P)$ is $P\otimes P$.
The signed shuffle product of the original odd elements is
\[
 x_u x_v=[y_u\mid y_v]-[y_v\mid y_u].
\]
Consequently any multiplicative simplicial extension would give a
simplicial linear map
\[
 h:P\otimes P\longrightarrow\Lambda^2P,
 \qquad h(u\otimes v-v\otimes u)=u\wedge v.
\]
It would retract the alternating inclusion
$\mathrm{alt}:\Lambda^2P\hookrightarrow P\otimes P$.
Eilenberg--Zilber gives $H_*(P\otimes P)=k[4]$, whereas
$H_3(\Lambda^2P)=k/2k\neq0$. A retraction would inject that
nonzero group into $H_3(P\otimes P)=0$, a contradiction.

Here is a complete finite verification of the latter nonzero group.
Use cut-set basis labels for $K(\Z,2)$. A pair in either square
functor is degenerate precisely when its two cut sets miss a common
position. The normalized exterior square has only two nonzero
terms, both $\Z^3$, in degrees three and four. Take bases
\[
 \begin{array}{ll}
 3:&12\wedge13,\quad12\wedge23,\quad13\wedge23,\\
 4:&12\wedge34,\quad13\wedge24,\quad14\wedge23.
 \end{array}
\]
Alternating vertex deletion gives
\[
 d_4=\begin{pmatrix}
 0&-1&1\\
 1&1&0\\
 0&-1&-1
 \end{pmatrix}.
\]
Its determinant is $-2$, and a two-by-two minor is a unit, so its
Smith invariants are $1,1,2$. Thus its third homology is $\Z/2$,
and localization at two gives the asserted $k/2k$.
The independently implemented finite certificate
\texttt{C\_bar\_boundary\_countercheck.py} verifies all relevant
faces, the normalized chain identities, this determinantal-divisor
calculation, and the tensor-square vanishing.
\end{proof}

The prescribed weight in this counterexample is below the eventual
row bound. The result therefore does not contradict the finite-range
extension theorems, and it is not a counterexample to C.3, which is
proved over $\Z_{(2)}$. It rules out an automatic bar-boundary map
based only on vertical exactness and eventual admissibility. Delayed
closures and different common-source constructions remain possible.
